\notepage{46}{The endpoints of an angular-momentum ladder}
For a fixed $\alpha$, successive applications of the ladder operators
change $\beta$ in steps of $\hbar$:
\[
 \cdots,\ \beta-2\hbar,\ \beta-\hbar,\ \beta,\
 \beta+\hbar,\ \beta+2\hbar,\ \cdots.
\]
Classically, $L_z^2\le L^2$. Quantum mechanically, for a normalized state,
\begin{align*}
 \alpha-\beta^2
 &=\mel{\alpha,\beta}{L^2-L_z^2}{\alpha,\beta}\\
 &=\mel{\alpha,\beta}{L_x^\dagger L_x}{\alpha,\beta}
  +\mel{\alpha,\beta}{L_y^\dagger L_y}{\alpha,\beta}\ge0.
\end{align*}
Hence $\alpha\ge0$ and $-\sqrt\alpha\le\beta\le\sqrt\alpha$.
The ladder must terminate above and below:
\[
 L_+\ket{\alpha,\beta_{\max}}=0,\qquad
 L_-\ket{\alpha,\beta_{\min}}=0.
\]
At the upper endpoint,
\begin{align*}
 0&=\mel{\alpha,\beta_{\max}}{L_-L_+}{\alpha,\beta_{\max}}\\
  &=\alpha-\beta_{\max}^2-\hbar\beta_{\max},
\end{align*}
so $\alpha=\beta_{\max}(\beta_{\max}+\hbar)$.
Similarly, using $L_+L_-=L^2-L_z^2+\hbar L_z$,
\[
 \alpha=\beta_{\min}(\beta_{\min}-\hbar).
\]
Equating these expressions and using $\beta_{\max}\ge\beta_{\min}$ yields
\[
 \beta_{\min}=-\beta_{\max},\qquad\beta_{\max}\ge0.
\]
There is a nonnegative integer $k$ such that
\[
 2\beta_{\max}=\beta_{\max}-\beta_{\min}=k\hbar.
\]
Writing $\beta_{\max}=\ell\hbar$ gives
\[
 \ell=0,\tfrac12,1,\tfrac32,2,\ldots,
\]
the nonnegative integer and half-integer possibilities from the
commutation algebra alone.

\begin{center}
\begin{tikzpicture}[>=Stealth,scale=0.85]
 \draw[->] (0,0)--(8,0) node[right]{$\beta$};
 \foreach \x/\t in {0.6/{-\ell\hbar},2.3/{\cdots},4/{m\hbar},5.7/{\cdots},7.4/{\ell\hbar}}{
  \fill (\x,0) circle (2pt);\node[below] at (\x,-0.08){$\t$};}
 \draw[->] (3.1,0.5)--(4.8,0.5) node[midway,above]{$L_+:\ +\hbar$};
\end{tikzpicture}
\end{center}
